\scanpage{58}
\mspage{3}
% archive foliation: 60
et telle que

a) $f\mid(\operatorname{Int}H)\cap\mathbb D^\circ$ est injective, son image $\subset\mathbb C\setminus K$

b) $f(0)=x$, \deleted{$\unsure{f(\pm t)=f(t)}$}

c) $f\mid]-1,0]$ et $f\mid[0,1[$ injectives, avec
même image $\subset K$.

\uline{Scholie} Chaque rive de $K$ \deleted{\unsure{est un}}\inserted{munit} $K$ d'une
structure analytique réelle canonique, via
l'inclusion de $K$ dans $\mathbb C$ (ou dans une
surface conforme quelc), mais le passage
de l'une à l'autre se fait par une
application (identique) qui n'est pas un
iso analytique : les propriétés de régularité
(caractère $C^\infty$, ou analytique réel etc)
reflètent sans doute les propriétés de
régularité correspondantes de $K$ dans $\mathbb C$.
(Ceci \inserted{devrait} se \unsure{magnifier|préciser} dans \deleted{\unsure{le th}}\inserted{un} Scholie
plus circonstancié au Théorème suivant).

\textcircled{3} \uline{Recollement}. On doit avoir un
théorème \unsure{dans le} style du
